%% notes-tex/modeling/scaling-laws/sections/7-sources.tex

\section{Sources\secstatus{todo}}
\label{sec:sources}

None of the five works below is in the mirror, so each is named by arXiv
identifier and flagged where its numbers are used. Adding them to the library is
a curation decision that has not been made; once it is, the citations here become
\texttt{\textbackslash cite} keys and the flags come off.

\begin{itemize}
\item Kaplan et al., 2020, \emph{Scaling Laws for Neural Language Models},
arXiv:2001.08361. The single-axis forms of Equation~\ref{eq:kaplan}, the
non-embedding parameter count, and the $6ND$ rule.
\item Hoffmann et al., 2022, \emph{Training Compute-Optimal Large Language
Models}, arXiv:2203.15556. The joint form with a floor, the three fitting
approaches (its appendix covers all three, including the Huber residual), the
IsoFLOP construction, and the Approach~3 constants.
\item Besiroglu et al., 2024, \emph{Chinchilla Scaling: A Replication Attempt},
arXiv:2404.10102. The refit of the Approach~3 data and the second set of
constants in Table~\ref{tab:constants}.
\item Yang et al., 2022, \emph{Tensor Programs V: Tuning Large Neural Networks
via Zero-Shot Hyperparameter Transfer}, arXiv:2203.03466. $\mu$P, the
parametrization under which the learning-rate optimum transfers across widths.
\item Muennighoff et al., 2023, \emph{Scaling Data-Constrained Language Models},
arXiv:2305.16264. The effective-data form for repeated epochs.
\end{itemize}

Within the repository: the 025 additive-baselines document
(\file{notes-tex/multimodal/025-additive-baselines/}) for the record count, the
splits and the leakage measurement quoted in Section~\ref{sec:counting}, and
\file{notes/torchcell.datasets.scerevisiae.costanzo2016.noise-computation.md} for
the replicate structure behind the floor estimate.
